% Source: analysis/defense-evaluation-20260928/rq4-table-compact-preview/build_preview.py.
% Reviewed layout; pooled canonical successes / COMPLETED, not mean subgroup ASR.
\begin{table}[!t]
\centering
\begin{minipage}{0.58\linewidth}
\centering
\setlength{\belowcaptionskip}{4pt}
\caption{SRDA attack success rates (ASR, \%) under defenses.}
\label{tab:defense-asr}
\small
\setlength{\tabcolsep}{4pt}
\renewcommand{\arraystretch}{1.02}
\begin{tabular}{@{}lrrrr@{}}
\toprule
Defense & \(G(1,3)\) & \(G(3,3)\) & \(G^{+}(3,3)\) & Overall \\
\midrule
None & 83.95 & 54.76 & 92.10 & 76.95 \\
ProtectAI & 74.57 & 44.44 & 89.14 & 69.38 \\
Prompt Guard & 21.48 & 7.65 & 26.54 & 18.56 \\
LLM Detector & 81.31 & 51.36 & 91.39 & 75.06 \\
\bottomrule
\end{tabular}
\par\smallskip
\begin{minipage}{\linewidth}
\footnotesize
\(G^{+}(3,3)\) denotes \(G(3,3)\) with two additional target-reaching links.
Pooled over six models and three paired seeds, with handoff guidance and
Natural feedback enabled. ASR uses completed executions
(810 scheduled per topology--defense condition).
\end{minipage}
\end{minipage}
\end{table}
